\documentclass[11pt]{article}
\usepackage[margin=1in]{geometry}
\usepackage[T1]{fontenc}
\usepackage[utf8]{inputenc}
\usepackage{array}
\usepackage{booktabs}
\usepackage{enumitem}
\usepackage{url}
\usepackage{xcolor}
\usepackage{listings}
\def\UrlBreaks{\do\/\do-\do\_\do\.\do\?\do\=\do\&}

\lstdefinestyle{cmd}{
  basicstyle=\ttfamily\small,
  columns=fullflexible,
  breaklines=true,
  frame=single,
  backgroundcolor=\color{black!3}
}
\setlength{\emergencystretch}{3em}

\title{ES-FA / BTP Setup Notes}
\author{First local setup for the KV260-based demo}
\date{Updated \today}

\begin{document}
\maketitle

\section*{Before starting}
These notes are written as a practical setup record for the ES-FA / BTP project,
not as a paper. The expected use case is simple: prepare the KV260 boot media,
set up the local machine once, pull the \texttt{btp} repository, run the project,
and generate the reports/plots/PDFs inside the local repo.

No project step here assumes cloud execution. Internet access is only needed to
download official installers/images or to clone/pull the GitHub repository. All
training outputs, Vivado/xsim logs, plots, runtime logs, and PDFs should stay
under local repo folders such as \path{results/}, \path{paper_output/}, and
\path{docs/}.

\section{1. Prepare the KV260 SD Card First}
\subsection{Image to flash}
Start with the board, because the rest of the setup is easier to understand once
the target platform is clear. The KV260 has QSPI boot firmware on the SOM and
uses the microSD card as the secondary boot device for Linux. For this project,
flash an official Kria Ubuntu image onto the microSD card.

\begin{itemize}[leftmargin=*]
\item First choice for a fresh setup: \textbf{Kria Ubuntu 24.04 LTS}.
\item Use Ubuntu 22.04 only if a specific Kria app or tutorial you need still
targets 22.04.
\item For older boards, check boot firmware compatibility before blaming the SD
card or project code. Newer Ubuntu images may require newer K26 boot firmware.
\end{itemize}

Download page:
\begin{lstlisting}[style=cmd]
https://ubuntu.com/download/amd-xilinx
\end{lstlisting}

\subsection{Flashing steps}
\begin{enumerate}[leftmargin=*]
\item Download the Kria Ubuntu image for KV260/K26.
\item Install Raspberry Pi Imager or Win32 Disk Imager on the host machine.
\item Insert a 32 GB or larger microSD card.
\item Select the downloaded image and write it to the card.
\item Eject the card safely after the write completes.
\end{enumerate}

\subsection{First boot checklist}
\begin{enumerate}[leftmargin=*]
\item Insert the flashed microSD card into the KV260 carrier.
\item Connect Ethernet and either HDMI/keyboard or USB-UART serial.
\item Power the board.
\item Login with the default Ubuntu account if prompted:
\begin{lstlisting}[style=cmd]
username: ubuntu
password: ubuntu
\end{lstlisting}
\item Change the password when Ubuntu asks.
\item Update the board:
\begin{lstlisting}[style=cmd]
sudo apt update
sudo apt upgrade
\end{lstlisting}
\item Keep the serial console available during the first few boots. It is the
fastest way to see SD-card or firmware boot problems.
\end{enumerate}

\section{2. Host Machine Requirements}
The repo has been run from a Windows host, and the scripts also have Bash
equivalents. For the presentation setup, the host machine should have:

\begin{itemize}[leftmargin=*]
\item Python 3.10, 3.11, or 3.12;
\item Git;
\item AMD Vivado/Vitis if hardware validation is required;
\item enough free disk space for \path{results/} and Vivado report folders.
\end{itemize}

Check the Xilinx tools locally:
\begin{lstlisting}[style=cmd]
vivado -version
xvlog -version
xelab -version
xsim -version
vitis -version
\end{lstlisting}

If these commands are not found, add the Xilinx \path{bin} folder to
\texttt{PATH}, or set one of the environment variables that the helper scripts
already search:
\begin{lstlisting}[style=cmd]
XILINX_VIVADO
XILINX_VITIS
XILINX_HOME
XILINX_INSTALL
VIVADO_HOME
VITIS_HOME
\end{lstlisting}

\section{3. Get the \texttt{btp} Repository}
The GitHub repository name is \texttt{btp}. Clone it like this:
\begin{lstlisting}[style=cmd]
git clone https://github.com/<your-github-user-or-org>/btp.git
cd btp
\end{lstlisting}

For an existing checkout:
\begin{lstlisting}[style=cmd]
cd btp
git status
git pull
\end{lstlisting}

If the folder was copied manually and is not a Git checkout, the code can still
run. For long-term work, clone a clean \texttt{btp} repo and copy only the files
you intentionally changed.

\section{4. Python Setup}
\subsection{Windows PowerShell}
\begin{lstlisting}[style=cmd]
cd btp
py -3.12 -m venv .venv312
.\.venv312\Scripts\Activate.ps1
python -m pip install --upgrade pip
python -m pip install -r requirements.txt
\end{lstlisting}

The helper script does the pip part:
\begin{lstlisting}[style=cmd]
.\scripts\setup_env.ps1 -PythonExe .\.venv312\Scripts\python.exe
\end{lstlisting}

\subsection{Linux or KV260 Ubuntu shell}
\begin{lstlisting}[style=cmd]
cd btp
python3 -m venv .venv312
source .venv312/bin/activate
python -m pip install --upgrade pip
python -m pip install -r requirements.txt
\end{lstlisting}

Helper script:
\begin{lstlisting}[style=cmd]
PYTHON_EXE=python ./scripts/setup_env.sh
\end{lstlisting}

\subsection{Quick sanity check}
\begin{lstlisting}[style=cmd]
python -c "import torch, torchvision, matplotlib; import paper1_training.training_core; print('imports ok')"
\end{lstlisting}

\section{5. Run the Local Demo}
\subsection{Fast smoke run}
For a presentation check, run the smoke pipeline first. This gives a quick
end-to-end run without waiting for full Vivado implementation:
\begin{lstlisting}[style=cmd]
.\scripts\run_esfa_pipeline.ps1 -PythonExe .\.venv312\Scripts\python.exe -Mode smoke -SkipVivado -SkipHardware
\end{lstlisting}

Bash equivalent:
\begin{lstlisting}[style=cmd]
MODE=smoke SKIP_VIVADO=1 SKIP_HARDWARE=1 ./scripts/run_esfa_pipeline.sh
\end{lstlisting}

\subsection{Manual software run}
Use this when you want to see the stages one by one:
\begin{lstlisting}[style=cmd]
python paper1_training/train_baseline.py
python experiments/run_first2.py
python experiments/run_remaining.py
python iterations/run_all.py
python analysis/compare.py
python analysis/evidence_report.py
\end{lstlisting}

\subsection{CLI demo}
These are the small commands I use to show that the project has a system layer,
not just training scripts:
\begin{lstlisting}[style=cmd]
python deployment/cli_interface/main.py --query "hi"
python deployment/cli_interface/main.py --query "strike rate 167.55 balls 39"
python deployment/cli_interface/main.py --query "classify results/runtime/sample_signal.bin"
\end{lstlisting}

\section{6. Hardware Validation}
After Vivado/xsim is available, run:
\begin{lstlisting}[style=cmd]
python hardware_validation/kv260/scripts/run_hw_validation.py --model-id baseline_paper1 --scheduler-mode both --clock-mhz 100
\end{lstlisting}

If you only want simulator-side validation:
\begin{lstlisting}[style=cmd]
python hardware_validation/kv260/scripts/run_hw_validation.py --model-id baseline_paper1 --scheduler-mode both --clock-mhz 100 --skip-vivado
\end{lstlisting}

The run folder is:
\begin{lstlisting}[style=cmd]
results/hardware_validation/<model-id>/<dense-or-event>/<run-id>/
\end{lstlisting}

The important files to check are \path{xsim_metrics.json},
\path{vivado_metrics.json}, \path{hardware_metrics.json}, and the logs/reports
inside the same run folder.

\section{7. Where Outputs Should Appear}
\begin{center}
\begin{tabular}{>{\ttfamily}p{0.34\linewidth}p{0.56\linewidth}}
\toprule
Folder & What to look for \\
\midrule
\path{results/baseline/} & baseline metrics, checkpoints, manifests, exports \\
\path{results/experiments/} & experiment metrics and exported artifacts \\
\path{results/plots/} & accuracy/energy/sparsity/latency/raster plots \\
\path{results/analysis/evidence/} & raw metric tables and percentage comparisons \\
\path{results/hardware_validation/} & xsim logs, Vivado reports, parsed hardware JSON \\
\path{results/runtime/} & CLI sample signal and adaptive decision logs \\
\path{paper_output/} & LaTeX report sources and compiled PDFs \\
\path{docs/} & setup notes, architecture notes, manual, checkpoints \\
\bottomrule
\end{tabular}
\end{center}

\section{8. Compile the PDFs Locally}
Install a local LaTeX distribution such as MiKTeX or TeX Live. Then compile from
the repository root:
\begin{lstlisting}[style=cmd]
pdflatex -interaction=nonstopmode -halt-on-error -output-directory=paper_output docs/first_time_user_setup_guide.tex
pdflatex -interaction=nonstopmode -halt-on-error -output-directory=paper_output docs/first_time_user_setup_guide.tex
pdflatex -interaction=nonstopmode -halt-on-error -output-directory=paper_output paper_output/progress_report.tex
pdflatex -interaction=nonstopmode -halt-on-error -output-directory=paper_output paper_output/progress_report.tex
\end{lstlisting}

Expected PDFs:
\begin{itemize}[leftmargin=*]
\item \path{paper_output/first_time_user_setup_guide.pdf}
\item \path{paper_output/progress_report.pdf}
\end{itemize}

\section{9. Common Problems}
\begin{center}
\begin{tabular}{p{0.38\linewidth}p{0.52\linewidth}}
\toprule
Problem & What I would check first \\
\midrule
\texttt{git} is not recognized & Install Git and reopen the terminal. \\
\texttt{ModuleNotFoundError: torch} & Activate \path{.venv312}, then rerun \texttt{python -m pip install -r requirements.txt}. \\
PyTorch install fails & Use Python 3.10--3.12, not a preview Python build. \\
\texttt{vivado} or \texttt{xsim} not found & Check \texttt{PATH} or set \texttt{XILINX\_VIVADO}/\texttt{XILINX\_VITIS}. \\
xsim compile fails & Open \path{results/hardware_validation/.../logs/xvlog_*.log}. \\
Negative Vivado slack & This is a design/timing result, not a failed setup. Save it and use it for timing closure. \\
KV260 does not boot & Reflash the SD card, check serial output, and confirm boot firmware compatibility. \\
CLI cannot find signal file & Use a path relative to repo root, for example \path{results/runtime/sample_signal.bin}. \\
No plots are produced & Rerun \texttt{python analysis/compare.py} after checking that metric JSON files exist. \\
\bottomrule
\end{tabular}
\end{center}

\section{10. Presentation Checklist}
\begin{enumerate}[leftmargin=*]
\item KV260 SD card is flashed and the board boots Ubuntu.
\item Python environment is active and the import check passes.
\item Smoke pipeline completes locally.
\item \path{results/ranking_table.md} and \path{results/analysis_summary.json}
exist.
\item Plots exist under \path{results/plots/}.
\item CLI examples return JSON output.
\item Hardware validation output exists if Vivado/xsim was available.
\item The setup guide and progress report PDFs are in \path{paper_output/}.
\end{enumerate}

\section{References Used for Board Setup}
\begingroup
\small
\begin{itemize}[leftmargin=*]
\item AMD KV260 Ubuntu SD-card setup:
\url{https://xilinx.github.io/kria-apps-docs/kv260/2022.1/linux_boot/ubuntu_24_04/build/html/docs/sdcard.html}
\item Ubuntu AMD-Xilinx image download:
\url{https://ubuntu.com/download/amd-xilinx}
\item AMD Kria SOM and Starter Kit image/firmware matrix:
\url{https://xilinx-wiki.atlassian.net/wiki/pages/viewpage.action?pageId=3307995140}
\item AMD KV260 boot device overview:
\url{https://docs.amd.com/r/en-US/ug1089-kv260-starter-kit/Boot-Devices-and-Firmware-Overview}
\end{itemize}
\endgroup

\end{document}
